\begin{resumo}

Este trabalho analisa vulnerabilidades críticas em sistemas web baseadas no OWASP Top 10 \cite{None2008RobustDefensesCrossSiteRequestForgery}, utilizando a defesa em profundidade como princípio unificador para integrar ferramentas de análise (SAST, DAST, IAST) e proteção em \textit{runtime} (RASP) \cite{Shahriar2011TaxonomyClassificationAutomaticMonitoringSecurityVulnerability, DeJesusDominguezGarcia2023SecurityTestingWebApplicationsSystematicLiteratureReview, Barth2022SystematicApproachWebProtectionRuntimeTools}. A metodologia consistiu em levantamento bibliográfico via \textit{Google Scholar} com suporte de IA generativa para síntese conceitual \cite{DeJesusDominguezGarcia2023SecurityTestingWebApplicationsSystematicLiteratureReview}. Identificaram-se lacunas na literatura sobre a segurança de \textit{shadow} APIs \cite{Escape2023StateGraphQLSecurity, Tanveer2025TowardsSecureAPIsRESTfulAPIVulnerabilityDetection} e a interpretabilidade de IA na detecção de ameaças \cite{Yadav2025CuttingEdgeApproachesIntrusionDetectionSystems, Alsaedi2023DeepLearningTechniqueEnabledWebApplicationFirewall}. Conclui-se que o futuro da área reside na automação via \textit{DevSecOps} \cite{Edmundson2022SANSDevSecOpsSurveyCreatingCulture} e na consolidação de arquiteturas \textit{Zero Trust} \cite{NIST2020SpecialPublication800207ZeroTrustArchitecture, Yadav2025ZeroTrustSecurityArchitectureMicroservicesMTLSOAuth20}.

\noindent \textbf{Palavras-chave:} Segurança Web, Defesa em Profundidade, \textit{DevSecOps}, \textit{Zero Trust}.

\end{resumo}

\begin{abstract}

This study analyzes critical web vulnerabilities based on the OWASP Top 10 \cite{None2008RobustDefensesCrossSiteRequestForgery}, using defense-in-depth as a unifying principle to integrate analysis tools (SAST, DAST, IAST) and runtime protection (RASP) \cite{Shahriar2011TaxonomyClassificationAutomaticMonitoringSecurityVulnerability, DeJesusDominguezGarcia2023SecurityTestingWebApplicationsSystematicLiteratureReview, Barth2022SystematicApproachWebProtectionRuntimeTools}. The methodology involved a literature review via \textit{Google Scholar} supported by generative AI for conceptual synthesis \cite{DeJesusDominguezGarcia2023SecurityTestingWebApplicationsSystematicLiteratureReview}. Gaps were identified regarding shadow API security \cite{Escape2023StateGraphQLSecurity, Tanveer2025TowardsSecureAPIsRESTfulAPIVulnerabilityDetection} and the interpretability of AI in threat detection \cite{Yadav2025CuttingEdgeApproachesIntrusionDetectionSystems, Alsaedi2023DeepLearningTechniqueEnabledWebApplicationFirewall}. The study concludes that the field's future lies in \textit{DevSecOps} automation \cite{Edmundson2022SANSDevSecOpsSurveyCreatingCulture} and the consolidation of \textit{Zero Trust} architectures \cite{NIST2020SpecialPublication800207ZeroTrustArchitecture, Yadav2025ZeroTrustSecurityArchitectureMicroservicesMTLSOAuth20}.

\noindent \textbf{Keywords:} Web Security, Defense-in-Depth, \textit{DevSecOps}, \textit{Zero Trust}.

\end{abstract}